\documentclass[11pt]{article}

\usepackage[margin=1in]{geometry}
\usepackage[utf8]{inputenc}
\usepackage[T1]{fontenc}
\usepackage{lmodern}
\usepackage{amsmath,amssymb,amsthm}
\usepackage{booktabs}
\usepackage{microtype}
\usepackage[numbers,sort&compress]{natbib}
\usepackage{xcolor}
\usepackage[colorlinks=true,linkcolor=blue!50!black,citecolor=blue!50!black,urlcolor=blue!50!black]{hyperref}

\newcommand{\draftnote}[1]{\textcolor{red!70!black}{[\textbf{Draft:} #1]}}

\title{Crystallising a Snapped Recurrent State:\\From an Annealed Snap to a Continuous, Quasi-Static Profile\\[0.5em]
\large Research note, 22 October 2026: not for circulation}

\author{
  David Barkhuizen\\
  Independent researcher
  \and
  Claude (Opus 5.5)\\
  Anthropic
}

\date{}

\begin{document}

\maketitle

\draftnote{A stand-alone note on the development of option~1 of the state-formation work, the
annealed crystallisation, from 20 October to the early morning of 22 October 2026, written at the
first author's request. It follows the status note of 21 October
(\texttt{state-formation-status.tex}), which it extends and in part supersedes. Its numbers come from
\texttt{indrajala-ml}'s recorded findings (\texttt{docs/state-formation-workplan.md}, decisions
D1--D59) at the commits in Section~\ref{sec:provenance}. Several registered probes were running at
its writing; their partial results are marked \emph{in flight} and are not verdicts.}

\begin{abstract}
Our method learns a finite-state transducer from input--output examples, with a recurrent state
snapped to learned codes, and certifies the learned machine exactly. Its main obstacle has been
forming that state and keeping it when the codes take it over. We describe how one remedy from the
literature, annealing the snap, became a working method in two days through a sequence of probes
each registered before its data. The milestones: a gradual handover from the continuous state to the
codes gave the first certificate (move-to-front on 4 items); a seeded crucible with a quasi-static
cool certified 8 of 8 runs of the word problem of $S_3$ (Wilson 95\% interval 0.68--1.00); and a
\emph{continuous} profile, in which every switch of that schedule became a smooth function of the
network's own readings, kept the rate on its design seeds (6 of 8), gave the crucible's first
certificate of move-to-front, and ran ten to fifty times faster. A probe with the coupling held at
zero certified 4 of 8 and refutes the claim that the handover is dispensable. Read in flight on fresh
seeds and new backdrops, the method finds the exact machine on nearly every run of $S_3$ and
move-to-front, and certifies runs of $S_4$ and a 12-state random machine; what fails is the stop,
which fires on the first right answers while some near ties still branch to wrong states. A
target-free reading of that failure (a tie whose branches split the machine's own minimised states)
separates the 19 networks examined without exception, and a stop that waits on it is in test. No
commitment term, traces or given decomposition were used.
\end{abstract}

\section{Setting}

The network (\texttt{snapped-window-recurrence.tex}) reads a window of input tokens and its current
state with a small causal transformer; the state is factored into $G$ parts, each part's next value
snapped to the nearest of its learned codes, so a trained network is a finite-state transducer. A
\emph{certificate} checks it against the target from every reachable state on every input: a proof,
not a sample. At the default snap ratio a code within the ratio of the nearest (a \emph{near tie})
branches, and both branches must be right; at the strict ratio 1.0 only the snap itself is checked.
A run \emph{succeeds} when the network it stops on certifies at both (D41). Every cell is feasible
(enough codes for the target's states; training reaches the exclusion bound $m + q - 1$; the
crucible's lengths reach every state and tell every pair apart), and every claim below was
registered before its data.

The obstacle (\texttt{state-formation-status.tex}): a state that never forms (\emph{formation}), or
forms and is lost when the discrete codes take over (\emph{conversion}). Of four options from the
literature review, option~1 anneals the snap: soft assignments hold every code at the data's mean
above a critical temperature and split as it cools
\citep{rose-1998-deterministic-annealing-clustering}, and quantising before a direction is learned
freezes it \citep{zhao-2026-continuous-first-discrete-later-vq-vae}. Soft-to-hard schedules for
vector quantisation \citep{agustsson-2017-soft-to-hard-vector-quantization} and annealed clustering
of a recurrent network's states into a machine
\citep{das-1993-unified-gradient-clustering-fsm-induction} are its nearest precedents.

\section{Development}

Table~\ref{tab:milestones} lists the milestones; the phases below give the reasoning between them.

\begin{table}[t]
\centering
\small
\begin{tabular}{@{}llll@{}}
\toprule
when & step & what changed & result \\
\midrule
20 Oct, morning & stage 4 & the annealed snap (D2--D12) & void: the anneal never ran \\
20 Oct & 4a, 4a$'$ & heat, hold, cool (D15--D24) & mechanics read; no certificate \\
20 Oct, 16:00 & 4a$''$ & the gradual handover (D25) & first certificate: mtf-4 seed 3 \\
20 Oct, evening & 4a$'''$ & the seeded crucible (D26, D27) & $S_3$: 2 of 8 \\
20 Oct, 20:30 & 4a$''''$ & the quasi-static cool (D28--D30) & $S_3$: 8 of 8 \\
21 Oct & D33 & the adaptive cool & $S_3$: 4 of 8 (its bar) \\
21 Oct & D35, D41--D43 & the answers-only stop; success & rules for discovery \\
21--22 Oct & D38 & the coupling held at 0 & $S_3$: 4 of 8, refuted \\
21 Oct, evening & D54 & the continuous profile (D48--D53) & $S_3$ 6 of 8; mtf-4 1 of 4 \\
21--22 Oct & D55--D58 & fresh seeds; $S_4$; random machines & in flight (Section~\ref{sec:now}) \\
22 Oct, 04:15 & D59 & the margin stop & in test \\
\bottomrule
\end{tabular}
\caption{Option~1's milestones. Counts are runs succeeding over runs; mtf-4 is move-to-front on 4
items.}
\label{tab:milestones}
\end{table}

\paragraph{The annealed snap (20 October, morning).} The first plan joined a continuous warm-up (the
codes bypassed, the cell's output the next state; D4), k-means codes from the warm-up's outputs (D6),
and a temperature halving until the codes were hard (D8), judged by a twelve-run panel. Its first
panel was void: the anneal never ran. The lesson was that a schedule needs a profile it is seen to
follow.

\paragraph{A thermal profile and the handover (20 October).} The first author's profile (D15): heat
to a full melt at twice the critical temperature (D16), hold, then cool exponentially, with thermal
noise (D17) and no fork (D18). The melting curve (D21) reads, on a run's own states, where each
factor melts and where it last splits. The cool's shapes and dropping the curriculum (D22--D24) said
little alone. The step that worked was the \emph{gradual handover} (D25): the next state is
$(1-l)\,h + l\,q$, the cell's output $h$ blended with its snapped code $q$, the coupling $l$ rising
from the melting point to the last split. Under it, move-to-front seed 3 certified: the method's
first certificate.

\paragraph{The crucible and quasi-staticity (20--21 October).} Groups such as $S_3$ have no
synchronising word and needed their state \emph{seeded} on short words first, so the crucible (D26)
holds the curriculum's warm-up material through the whole profile, then grows to longer lengths; a
population of seeds lets the certificate select (D27). It certified 2 of 8. Read on the six failures,
each lost a formed state where the cool or the coupling outran the network, so three fixes followed
(D28--D30): the warm-up runs to its cap; the cool's clock holds while the screen is below the hold's
best (the \emph{gate}: the cool waits for the network); and the coupling's window spans at least three
octaves. The quasi-static crucible certified 8 of 8 (Wilson 95\% interval 0.68--1.00).

\paragraph{Questioning it (21 October).} An adaptive cool paced by the codes' heat capacity
\citep{huang-1986-efficient-cooling-schedule} (D32, D33) met its bar exactly, 4 of 8 (0.22--0.78): its
steps outran the handover below the cascade, where the coupling rose by up to 0.65 in an epoch. Two
reviews tightened the rules: the run now stops on its answers alone, at the first step whose snapped
machine answers right with at most the state bound's states, the certificate only judging (D35, the
target out of the loop as discovery needs); a success is a certificate at both verdicts (D41); the
bound is exact (D43). The coupling probe (D38) asked whether the handover was needed at all, the codes
annealing on the continuous state but never entering it: 4 of 8, refuted. Its four failures never
once answered right on their snapped machine.

\paragraph{The continuous profile (21 October, evening).} Reading the coupling probe's seed 3 showed
the schedule's switches (a clock that holds or runs, a coupling window, a hard snap, growth in
blocks) jolting a network that should descend its energy profile smoothly. At the first author's
direction (``make everything we can, that still makes sense to do so, continuous. all at once''),
every switch became a smooth function of the network's readings (D48):
\begin{itemize}
  \item \emph{the clock} advances at a rate $g = \mathrm{clip}((s/\ell - 0.95)/0.04)$, the screen
  $s$ against the hold's level $\ell$, in one exponential cool halving every 30 epochs (D49);
  \item \emph{the coupling} is read live from each factor's sharpness $S$, its codes' mean largest
  soft weight: $l = \mathrm{clip}((S - S_0)/(1 - S_0))$ from the cool's start (D50);
  \item \emph{no hard switch}: soft to the end, with a straight-through gradient
  \citep{bengio-2013-straight-through-estimator} built into the Rust crate, each temperature held once
  its codes are hard to 0.1\% (D51);
  \item \emph{growth} in $\ln L$ at $0.02\,g\,l$ up to the cap, each epoch's lengths a mixture of
  the two nearest ranges (D52);
  \item \emph{every reading every epoch}: both certificates' breakdowns, the snapped answers, the
  melting curve every 25 epochs (D53).
\end{itemize}
A self-annealed stochastic quantiser \citep{takida-2022-sq-vae-self-annealed-stochastic-quantization}
is the nearest precedent for a schedule read from the codes themselves that we found. On the
quasi-static crucible's design seeds (D54) it certified $S_3$ 6 of 8 (0.41--0.93), two still
running, including seeds the adaptive cool lost; it gave the crucible's first certificate of
move-to-front (seed 0); its other three move-to-front runs stopped on machines exact at 1.0. Runs
took 11--57 minutes on $S_3$ and about an hour on move-to-front, against hours to days before.

\section{Where it stands}
\label{sec:now}

\paragraph{Fresh seeds and new backdrops (in flight).} The continuous cell, unchanged, on seeds no
run had touched (D55) and on three new backdrops (D56--D58), each registered before its data.
Table~\ref{tab:now} gives the counts at about 04:00 on 22 October.

\begin{table}[t]
\centering
\small
\begin{tabular}{@{}lrrrrl@{}}
\toprule
cell & succeeded & exact at 1.0 only & budget spent & running & bar \\
\midrule
$S_3$, seeds 24--39 (D55) & 7 & 5 & 2 & 2 & 12 of 16 \\
mtf-4, seeds 24--31 (D55) & 3 & 4 & 0 & 1 & 2 of 8 \\
$S_4$, seeds 0--2 (D57) & 1 & 0 & 0 & 2 & 1 of 3 \\
random, 12 states, seeds 0--3 (D58) & 1 & 0 & 0 & 3 & 2 of 4 \\
random, 16 states, seeds 0--3 (D56) & 0 & 0 & 3 & 1 & 1 of 4 \\
\bottomrule
\end{tabular}
\caption{The continuous cell on fresh seeds and new backdrops, in flight. ``Exact at 1.0 only'':
stopped by the answers-only stop on a network certified at the strict ratio but not the default.}
\label{tab:now}
\end{table}

Move-to-front has met its bar on fresh seeds, and $S_4$'s seed 1, which certified before only with
the commitment term, certified without it. $S_3$ will most likely miss its registered bar: 9 is the
most it can reach. Yet of the 14 $S_3$ runs finished, 12 found the exact machine, as did all 7 on
move-to-front. The 16-state random machine, on which every earlier arm failed, remains unsolved.

\paragraph{The margin.} Every run stopped exact at 1.0 alone failed the default verdict only through
its near ties: 1--13 transitions whose tie branches to a wrong state, among 2--188 near ties. Near ties
alone do not decide it: the ten runs that succeeded had 2--125. What decides it is whether a tie
\emph{splits} the machine: whether its branch leads to another state of the network's own minimised
machine \citep{moore-1956-gedanken-experiments-sequential-machines} than the snap's. That is the machine's own nondeterminism, read
without the target. On the 19 stopped networks of D55, the count of split ties was above zero on each
of the 9 early stops and zero on each of the 10 that succeeded. The answers-only stop fires at the
first right answers, before training has widened these margins; the \emph{margin stop} (D59) waits
until no tie splits, at most 200 steps past the first pass. Its probe, the 12 early stops re-run from
their seeds, was launched at 04:15 on 22 October; it counts machines lost during the wait as well as
gained.

\section{What is not established}
\label{sec:open}

\begin{itemize}
  \item \emph{The margin stop}: whether waiting turns exact machines into certificates, or loses
  them; the split-tie reading is shown on 19 networks of two machines.
  \item \emph{The rate}: the fresh-seed matrix's verdict is pending, and on $S_3$ is likely below its
  bar under the stop it ran with.
  \item \emph{The unstructured machine}: the 16-state random machine has no certificate in any arm;
  its failures are to be classified (extra codes against near ties).
  \item \emph{The large tier}: $S_5$, move-to-front on 5 items and division modulo 13 (120--169
  states) are untried under the continuous profile; each epoch there costs minutes and gigabytes.
  \item \emph{Design seeds}: the profile's constants (0.95--0.99, 30 epochs, 0.02, 0.999) were set
  on the design seeds and not tuned since; their sensitivity is unmeasured.
  \item \emph{Discovery proper} still needs a state bound; every rate here is at the exact bound.
\end{itemize}

\section{Provenance}
\label{sec:provenance}

All in \texttt{davidbarkhuizen/indrajala-ml}; the decisions and findings are in its
state-formation workplan (\texttt{docs/state-formation-workplan.md}). Up to D34 as in \texttt{state-formation-status.tex}. The
continuous profile's decisions (D47--D54) in 229bf56 (PR \#906); the straight-through gradient in
\texttt{indrajala-math-rust} 45d6a37 (PR \#70), bumped in fb0f1be (\#907); \texttt{--continuous} in
5368e0e (\#908), the continuous probe run at its tree. The fresh-seed matrix (D55) in 43acf1d (\#909);
the random and $S_4$ probes (D56--D58) in e934df2, dc351d5 and bc37b0b (\#910--\#912); the margin
stop in 61fcf3e (\#913), its probe in 7f7f09c (\#914). Every probe's script, with its expectations
written before its data and its SHA-256 recorded in the workplan, is \texttt{\textasciitilde/code/probe-d\textit{NN}.sh}.
Bit parity of the continuous cell checked across four machines (255 log rows identical on each).

\bibliographystyle{plainnat}
\bibliography{library,extra}

\end{document}
